\documentclass[10pt,a4paper]{amsart}
\usepackage{amsmath,amssymb,amsthm,graphicx}
\usepackage{verbatim}
\newenvironment{reference}{\comment}{\endcomment}
%\newenvironment{reference}{}{}
\newenvironment{slogan}{\comment}{\endcomment}
\newenvironment{history}{\comment}{\endcomment}

% For commutative diagrams we use Xy-pic
\usepackage[all]{xy}

% We use 2cell for 2-commutative diagrams.
\xyoption{2cell}
% Theorem environments.
%
\theoremstyle{plain}
\newtheorem{theorem}[subsection]{Theorem}
\newtheorem{proposition}[subsection]{Proposition}
\newtheorem{lemma}[subsection]{Lemma}

\theoremstyle{definition}
\newtheorem{definition}[subsection]{Definition}
\newtheorem{example}[subsection]{Example}
\newtheorem{exercise}[subsection]{Exercise}
\newtheorem{situation}[subsection]{Situation}

\theoremstyle{remark}
\newtheorem{remark}[subsection]{Remark}
\newtheorem{remarks}[subsection]{Remarks}

\numberwithin{equation}{subsection}

% Macros
%
\def\lim{\mathop{\mathrm{lim}}\nolimits}
\def\colim{\mathop{\mathrm{colim}}\nolimits}
\def\Spec{\mathop{\mathrm{Spec}}}
\def\Hom{\mathop{\mathrm{Hom}}\nolimits}
\def\Ext{\mathop{\mathrm{Ext}}\nolimits}
\def\SheafHom{\mathop{\mathcal{H}\!\mathit{om}}\nolimits}
\def\SheafExt{\mathop{\mathcal{E}\!\mathit{xt}}\nolimits}
\def\Sch{\mathit{Sch}}
\def\Mor{\mathop{\mathrm{Mor}}\nolimits}
\def\Ob{\mathop{\mathrm{Ob}}\nolimits}
\def\Sh{\mathop{\mathit{Sh}}\nolimits}
\def\NL{\mathop{N\!L}\nolimits}
\def\CH{\mathop{\mathrm{CH}}\nolimits}
\def\proetale{{pro\text{-}\acute{e}tale}}
\def\etale{{\acute{e}tale}}
\def\QCoh{\mathit{QCoh}}
\def\Ker{\mathop{\mathrm{Ker}}}
\def\Im{\mathop{\mathrm{Im}}}
\def\Coker{\mathop{\mathrm{Coker}}}
\def\Coim{\mathop{\mathrm{Coim}}}

% Boxtimes
%
\DeclareMathSymbol{\boxtimes}{\mathbin}{AMSa}{"02}

%
% Macros for moduli stacks/spaces
%
\def\QCohstack{\mathcal{QC}\!\mathit{oh}}
\def\Cohstack{\mathcal{C}\!\mathit{oh}}
\def\Spacesstack{\mathcal{S}\!\mathit{paces}}
\def\Quotfunctor{\mathrm{Quot}}
\def\Hilbfunctor{\mathrm{Hilb}}
\def\Curvesstack{\mathcal{C}\!\mathit{urves}}
\def\Polarizedstack{\mathcal{P}\!\mathit{olarized}}
\def\Complexesstack{\mathcal{C}\!\mathit{omplexes}}
% \Pic is the operator that assigns to X its picard group, usage \Pic(X)
% \Picardstack_{X/B} denotes the Picard stack of X over B
% \Picardfunctor_{X/B} denotes the Picard functor of X over B
\def\Pic{\mathop{\mathrm{Pic}}\nolimits}
\def\Picardstack{\mathcal{P}\!\mathit{ic}}
\def\Picardfunctor{\mathrm{Pic}}
\def\Deformationcategory{\mathcal{D}\!\mathit{ef}}
\usepackage{cleveref}
\crefname{originalequation}{original equation}{original equations}
\Crefname{originalequation}{Original equation}{Original equations}
\makeatletter
\newlabel{equation-C-prime}{{0AL4}{source}{}{sourcepointer-1}{}}
\newlabel{equation-C-prime@cref}{{[originalequation][0][]0AL4}{[1][0][]source}{}{}{}}
\makeatother
\crefname{originalreference}{original reference}{original references}
\Crefname{originalreference}{Original reference}{Original references}
\makeatletter
\newlabel{lemma-approximate-presentation-rig-smooth}{{0GAW}{source}{}{sourcepointer-2}{}}
\newlabel{lemma-approximate-presentation-rig-smooth@cref}{{[originalreference][0][]0GAW}{[1][0][]source}{}{}{}}
\makeatother
\crefname{originalreference}{original reference}{original references}
\Crefname{originalreference}{Original reference}{Original references}
\makeatletter
\newlabel{lemma-equivalent-with-artin}{{0AJU}{source}{}{sourcepointer-3}{}}
\newlabel{lemma-equivalent-with-artin@cref}{{[originalreference][0][]0AJU}{[1][0][]source}{}{}{}}
\makeatother
\crefname{originalreference}{original reference}{original references}
\Crefname{originalreference}{Original reference}{Original references}
\makeatletter
\newlabel{lemma-get-morphism-general-better}{{0GAQ}{source}{}{sourcepointer-4}{}}
\newlabel{lemma-get-morphism-general-better@cref}{{[originalreference][0][]0GAQ}{[1][0][]source}{}{}{}}
\makeatother
\crefname{originalreference}{original reference}{original references}
\Crefname{originalreference}{Original reference}{Original references}
\makeatletter
\newlabel{lemma-get-morphism-principal}{{0AK7}{source}{}{sourcepointer-5}{}}
\newlabel{lemma-get-morphism-principal@cref}{{[originalreference][0][]0AK7}{[1][0][]source}{}{}{}}
\makeatother
\title{Rig-etale algebras over a principal ideal are completions of finite type algebras}
\author{The Stacks Project Authors}
\date{stacks-project:restricted.tex:local-snapshot-20261012; GNU FDL 1.2 or later}
\begin{document}
\maketitle
\noindent\textit{Editorial note.} All delivered mathematical text is the author's own, copied from the pinned Stacks Project chapter file restricted.tex; no statement or proof step was rewritten, reordered or removed.\par\medskip
\noindent\textit{Editorial note.} The counted proof is the author's complete argument: the proof displays the concrete data (c, f\_1..f\_r, divisibility, inclusion) obtained from part (6) before the final invocation, so the argument is complete at the named-lemma level; the reference environment citing Elkik is statement-side metadata\par\medskip
\noindent\textit{Editorial note.} The label lemma-approximate-principal names the counted statement, and the proof environment immediately following it in the source is delivered verbatim as its proof.\par\medskip
\noindent\textit{Editorial note.} The counted statement is carried by the author's own lemma environment in the chapter source.\par\medskip
\noindent\textit{Allowed standard prerequisites (editorial).}\begin{itemize}
\item Lemma lemma-equivalent-with-artin part (6): for a presentation B = A[x\_1,...,x\_r]\textasciicircum{}/J over (A, I) it produces c >= 0 and f\_1,...,f\_r in J such that the Jacobian determinant det(d f\_j / d x\_i) divides a\textasciicircum{}c in B and a\textasciicircum{}c J is contained in (f\_1,...,f\_r) + J\textasciicircum{}2
\item Lemma lemma-approximate-presentation-rig-smooth (delivered, L48002): converts precisely such presentation data into a finite type A-algebra C with B isomorphic to C\textasciicircum{}
\end{itemize}\medskip
\ifdefined\hypertarget\hypertarget{sourcepointer-1}{}\fi
\noindent\textit{Source reference.} Original equation 0AL4: sources/stacks-restricted/pinned-original-source.tex, source lines 326--326. Prerequisite pointer invoked at 2286; it names a standard fact of the same chapter that the counted argument uses, and every step of the counted proof built on it is retained verbatim in the delivered body.\par\medskip
\ifdefined\hypertarget\hypertarget{sourcepointer-2}{}\fi
\noindent\textit{Source reference.} Original reference 0GAW: sources/stacks-restricted/pinned-original-source.tex, source lines 1638--1638. Prerequisite pointer invoked at 2299; it names a standard fact of the same chapter that the counted argument uses, and every step of the counted proof built on it is retained verbatim in the delivered body.\par\medskip
\ifdefined\hypertarget\hypertarget{sourcepointer-3}{}\fi
\noindent\textit{Source reference.} Original reference 0AJU: sources/stacks-restricted/pinned-original-source.tex, source lines 1810--1810. Prerequisite pointer invoked at 2294; it names a standard fact of the same chapter that the counted argument uses, and every step of the counted proof built on it is retained verbatim in the delivered body.\par\medskip
\ifdefined\hypertarget\hypertarget{sourcepointer-4}{}\fi
\noindent\textit{Source reference.} Original reference 0GAQ: sources/stacks-restricted/pinned-original-source.tex, source lines 1167--1167. Prerequisite pointer invoked at 2302; it names a standard fact of the same chapter that the counted argument uses, and every step of the counted proof built on it is retained verbatim in the delivered body.\par\medskip
\ifdefined\hypertarget\hypertarget{sourcepointer-5}{}\fi
\noindent\textit{Source reference.} Original reference 0AK7: sources/stacks-restricted/pinned-original-source.tex, source lines 1294--1294. Prerequisite pointer invoked at 2304; it names a standard fact of the same chapter that the counted argument uses, and every step of the counted proof built on it is retained verbatim in the delivered body.\par\medskip
\begin{lemma}
\label{lemma-approximate-principal}
\begin{reference}
The rig-\'etale case of \cite[III Theorem 7]{Elkik}
\end{reference}
Let $A$ be a Noetherian ring and $I = (a)$ a principal ideal.
Let $B$ be an object of (\ref{equation-C-prime}) which is
rig-\'etale over $(A, I)$.
Then there exists a finite type $A$-algebra $C$ and an
isomorphism $B \cong C^\wedge$.
\end{lemma}

\begin{proof}
Choose a presentation $B = A[x_1, \ldots, x_r]^\wedge/J$.
By Lemma \ref{lemma-equivalent-with-artin} part (6) we can find
$c \geq 0$ and $f_1, \ldots, f_r \in J$ such that
$\det_{1 \leq i, j \leq r}(\partial f_j/\partial x_i)$ divides
$a^c$ in $B$ and
$a^c J \subset (f_1, \ldots, f_r) + J^2$.
Hence Lemma \ref{lemma-approximate-presentation-rig-smooth} applies.
This finishes the proof, but we'd like to point out that in
this case the use of
Lemma \ref{lemma-get-morphism-general-better}
can be replaced by the much easier
Lemma \ref{lemma-get-morphism-principal}.
\end{proof}

\end{document}
